\scanpage{164}
\deleted{indécomposables distincts, soient $V_i$. pour $i,j$, $V_i\otimes V_j$ est un $\underline U$-module, donc se décompose en un nb fini de sous $\underline U$-modules \deleted{irréductibles} indécomposables $W_{ij\unsure{k|u}}$, l'ens des $W_{ijk}$ obtenus est fini. \deleted{Soit $\mathcal J'_V$ le noyau} de $U$ dans $\sum W_{ijk}$}

\deleted{Montrons qu'il existe un $W\in\inserted{K^{\underline U}}$ tel que, si $W_1$ et $W_2$ sont $\in K^{\underline U_V}$, alors $W_1\otimes W_2$, muni de sa structure ``naturelle'' de $\underline U$-module, ait un annulateur $\supset\mathcal J_W$. En effet, il suffit de le vérifier pour $W_1$ et $W_2$ indécomposables (\unsure{grâce à} \unsure{Rem th de|Krull} Schmidt) or il n'y a qu'un nb \uline{fini} de $V_i$ indécomposables distincts; on prendra $W=\sum V_i\otimes V_j$. Ceci dit, \deleted{on veut définir} \deleted{\unsure{D'ailleurs}, \illegible{2} un hom}}

Ceci est un cas particulier de la \deleted{\unsure{structure}}\inserted{donnée}
suivante. Tout bi-$\underline U$-module permis
(i.e. un module permis sur $\underline U\mathbin{\widehat\otimes}\underline U$)
définit un $\underline U$-module permis de
façon fonctorielle, i.e. on se donne
un \uuline{hom diagonal} $\underline U\longrightarrow\underline U\mathbin{\widehat\otimes}\underline U$. Mais on

% archive sheet number: 163
